\subsection*{Parameter derivatives and explicit terminal matching}
\label{subsec:A11-deriv-stop}
The previous source-native theorem controls the two relative
azimuthal bumps in (A.11). The dependence on the parameter $\eta$
is also tame, and the final constant-slope stopping length can be
determined without a numerical shooting argument.

\begin{proposition}[First-parameter-derivative bounds]
\label{prop:A11-eta-derivative}
Assume the exact Appendix-A stages and Theorem~\ref{thm:A11-true-corrections},
with $T_f=512$ and $0<\lambda\le1/20$. For the actual two correction
coefficients and the corrected azimuthal logarithmic slope,
\[
 \max_i\|\partial_\eta c_i\|_{L^\infty([-1,1])}
 <12\lambda^{28},\qquad
 \|\partial_\eta\ell_c\|_{L^\infty}<5\cdot10^6\lambda^{28}.
\]
The pressure-neutrality and angular-moment identities also
hold after differentiating in $\eta$, since the correction equations
hold identically and the bump integrals have compact support.
\end{proposition}
\begin{proof}
Write $r=I/(XH)$. It satisfies $r'+(1+\ell)r=1$ in logarithmic
radial coordinates. The reference value is $r=5/8$ and is
independent of $\eta$; all preceding unedited slopes are also
$\eta$-independent. During the interpolation (A.10),
\[
 \partial_\eta\ell
 =\frac{\sigma'(y/T_f)}{T_f}\frac{2\eta}{1+\eta^2},
 \quad |\partial_\eta\ell|\le64/T_f=1/8.
\]
Since $0\le r\le20/17$ and $1+\ell\ge17/20$, the
equation for $v=\partial_\eta r$ gives $|v|\le50/289$.
On the following $30\log(1/\lambda)-3$ unedited units,
$v$ decays at rate $1-\lambda$. Thus the moment discrepancy
$\Delta_I$ defined in \eqref{eq:A11-two-exact} obeys
\[
 |\partial_\eta\Delta_I|\le\frac{50}{289}
 e^{3(1-\lambda)}\lambda^{30(1-\lambda)}
 <\frac{1050}{289}\lambda^{28}<4\lambda^{28}.
\]
The exact nonlinear correction system is $B_{11}c+Q(c,c)
=(\Delta_I,0)^T$. By Theorem~\ref{thm:A11-matrix-bound}
and \eqref{eq:A11-c-coeff},
\[
 \|B_{11}^{-1}DQ(c)\|_\infty
 \le(3/2)(640)(192)\lambda^{28}<1/2.
\]
Differentiate the system and invert its Jacobian to obtain
$\|c_\eta\|_\infty\le3|\partial_\eta\Delta_I|
<12\lambda^{28}$.

Put $g=\sum c_i\beta_i$. On the two disjoint supports
$\ell_c=-\lambda+g_y/(1+g)$.
Using $|g|<1/2$, $\|\beta\|_\infty\le640/3$ and
$\|\beta'\|_\infty<200000$ yields
\[
 |\partial_\eta\ell_c|\le
 2(200000)(12)\lambda^{28}
 +4(200000)(192)(640/3)(12)\lambda^{56}
 <5\cdot10^6\lambda^{28}.
\]
\end{proof}

\begin{proposition}[Quantitative terminal stopping length]
\label{prop:A11-terminal-stop}
Assume the actual (A.11) correction and the subsequent
source-prescribed one-unit transition to $\ell=-1$, the hold
$\ell=-1$ for $4\log(1/h)$, and one-unit transition to
$\ell=-h$, with $h<\lambda\le1/20$ and $0<h\le1/100$.
For the terminal cutoff (A.12), fix $c_o=1/256$.
Let $Q_{\rm start}$ denote the value of the source $Q_s$
at the start of the last $\ell=-h$ hold, and $Q_p$ the
cutoff matching value (A.13). Then
\begin{equation}\label{eq:A11-terminal-Q-bounds}
 Q_{\rm start}\ge\frac4e(1-h)\log(1/h)>5,\qquad
 \frac{c_oh}{1-c_oh}\le Q_p
 \le\frac{e^{3(1-h)}c_oh}{1-c_oh}<\frac h{10}.
\end{equation}
In particular the unique positive holding length is
\begin{equation}\label{eq:A11-stop-exact}
 L_{\rm stop}=\frac1{1-h}\log\frac{Q_{\rm start}}{Q_p}>0.
\end{equation}
It is independent of $\eta$.
\end{proposition}
\begin{proof}
The restored angular moment implies
$Q_s=(\lambda-h)/(1-\lambda)>0$ at the exit of (A.11).
The source ODE is $Q_s'+(1+\ell)Q_s=-\ell-h$.
The first transition has nonnegative right-hand side;
on the $\ell=-1$ hold, $Q_s'=1-h$, so its exit value is
at least $4(1-h)\log(1/h)$. During the following unit transition
the right-hand side remains nonnegative and $0\le1+\ell\le1$,
giving the stated lower bound on $Q_{\rm start}$.
For the terminal cutoff $f_o=1-c_oh\psi_o$ of (A.12),
$f_o(0)=1-c_oh$, $f_o(3)=1$, and $f_o'\ge0$.
The integrating factor in (A.13) simplifies exactly to
\[
 Q_p=\frac1{f_o(0)}\int_0^3e^{(1-h)v}f_o'(v)\,dv.
\]
Integration of $f_o'$ gives the displayed interval for $Q_p$.
The numerical inequalities follow from $e^3<21$, $e<3$,
$\log100>4$ and $h\le1/100$. Finally on $\ell=-h$
the equation is $Q_s'=-(1-h)Q_s$; the unique positive length
at which $Q_s$ reaches $Q_p$ is \eqref{eq:A11-stop-exact}.
The exit moment and all later prescribed slopes are
$\eta$-independent, so the length is as well.
\end{proof}

These estimates quantify source constructions that were previously
given qualitatively. They do not certify all of the source's
stress-cone inequalities, the full moment-matched numerical exterior,
Appendix-B contraction or the unforced $n=1,2$ correction systems.
